\chapter{Qué se puede leer en un jacobiano}
\label{cap:usos}

Los jacobianos no son solo un paso intermedio para llegar a las funciones
impulso-respuesta. Son objetos económicos con contenido propio, y aprender a
leerlos es una de las habilidades más transferibles que da este método.

\section{Funciones impulso-respuesta y descomposiciones}

Lo básico ya está dicho: $d\mathbf{X} = \mathbf{G}\,d\mathbf{Z}$, y calcular una
respuesta adicional cuesta un producto matriz--vector. Dos usos menos obvios:

\paragraph{Descomponer una respuesta en canales} Supongamos que el consumo
responde a una política monetaria a través de tres canales: el efecto directo de
la tasa, el efecto ingreso vía salarios, y el efecto vía dividendos. Como la
respuesta total es
$d\mathbf{C} = \J^{C,r}d\mathbf{r} + \J^{C,w}d\mathbf{w} + \J^{C,d}d\mathbf{d}$,
basta con evaluar cada término por separado usando las respuestas de equilibrio
general $d\mathbf{r}$, $d\mathbf{w}$, $d\mathbf{d}$. La suma es exacta, no
aproximada, porque el modelo es lineal.

\paragraph{Contrafactuales de ``apagar un canal''} ¿Qué pasaría si los hogares no
respondieran a la tasa de interés? Se resuelve el equilibrio general reemplazando
$\J^{C,r}$ por cero y dejando todo lo demás igual. Esto es un experimento
coherente y barato, y es muchísimo más informativo que comparar dos
calibraciones distintas.

\section{Los jacobianos como estadísticos suficientes}

El resultado conceptual más importante de esta literatura es que, para una clase
amplia de preguntas, lo único que hace falta saber del lado de los hogares es un
jacobiano. Auclert, Rognlie y Straub (2018) muestran que en una clase de modelos
keynesianos la respuesta del producto a la política fiscal depende del modelo de
hogares \emph{solo} a través de $\J^{C,y}$, la matriz de propensiones marginales
a consumir intertemporales (las ``iMPC'').

Esto tiene dos consecuencias prácticas.

\begin{enumerate}[leftmargin=1.4em]
\item \textbf{Se puede disciplinar el modelo con microdatos.} La primera columna
de $\J^{C,y}$ es el perfil de respuesta del consumo a una transferencia
inesperada: exactamente lo que estiman los estudios de reembolsos tributarios y
loterías. Un modelo cuyo $\J^{C,y}$ no reproduce esos perfiles está mal
calibrado, aunque reproduzca la distribución de riqueza.
\item \textbf{Se pueden comparar modelos sin resolverlos.} Dos modelos con el
mismo $\J^{C,y}$ dan las mismas respuestas agregadas. Si un modelo con dos
activos y uno con un activo tienen jacobianos parecidos, la complejidad adicional
no está comprando nada para esa pregunta.
\end{enumerate}

\begin{idea}
Cuando leas un artículo con un modelo HANK complicado, la primera pregunta útil
no es ``¿cómo calibraron el proceso de ingreso?'' sino ``¿cómo se ve
$\J^{C,y}$?''. Si el artículo no lo muestra, pedirlo es la mejor pregunta que
puedes hacer en un seminario.
\end{idea}

\section{Leer la forma de un jacobiano}

Con algo de práctica, la forma de la matriz dice qué mecanismo domina. Un
catálogo breve:

\begin{itemize}[leftmargin=1.4em]
\item \textbf{Filas constantes} (como en \eqref{eq:jac-ri}): suavizamiento
perfecto, el agente solo mira valores presentes. Es el caso de ingreso
permanente.
\item \textbf{Diagonal dominante}: restricciones de liquidez. El agente responde
a lo que pasa hoy mucho más que a lo que pasa mañana, porque no puede trasladar
recursos. Cuanto más alta la fracción de restringidos, más dominante la diagonal.
\item \textbf{Columnas en forma de carpa}: sustitución intertemporal con
anticipación, como en la figura \ref{fig:columnas}.
\item \textbf{Triangular inferior}: no hay anticipación. Si esto aparece en un
bloque donde debería haber agentes prospectivos, es un error de programación.
\item \textbf{Primera fila igual a cero salvo en $s=0$}: los agregados de hoy
dependen solo de los insumos de hoy y de la distribución predeterminada. Es la
firma de una variable de acervo.
\end{itemize}

El último caso aparece en la parte V y conviene retenerlo, porque es una
verificación gratuita: si una variable agregada está determinada por el acervo
heredado, ninguna noticia sobre el futuro puede moverla hoy.

\section{La matriz de equilibrio general $\mathbf{G}$}

Después de resolver, $\mathbf{G}^{o,z}$ es una matriz $T\times T$ cuya columna
$s$ es la respuesta de equilibrio general de la variable $o$ a una noticia de que
el choque $z$ ocurrirá en $s$. Es útil pensarla como ``todas las funciones
impulso-respuesta del modelo, a la vez''.

Dos diferencias importantes respecto del jacobiano parcial $\J$:

\begin{itemize}[leftmargin=1.4em]
\item $\mathbf{G}$ incorpora la retroalimentación de precios. $\J^{C,r}$ dice
cómo responde el consumo si la tasa cambia; $\mathbf{G}^{C,Z}$ dice cómo responde
el consumo cuando la productividad cambia \emph{y} la tasa se ajusta en
equilibrio.
\item $\mathbf{G}$ depende de todo el modelo; $\J$ solo del bloque. Si cambia la
regla de política monetaria, $\mathbf{G}$ cambia y $\J$ no. Es exactamente la
distinción que la crítica de Lucas vuelve relevante, y el método la hace
operativa: los objetos invariantes a la política son los jacobianos de bloque.
\end{itemize}
